% Compile twice: lualatex 30000347.tex
\documentclass[11pt]{article}
\usepackage[margin=25mm]{geometry}
\usepackage{fontspec,unicode-math}
\setmainfont{Latin Modern Roman}
\setmathfont{Latin Modern Math}
\usepackage{amsmath,amssymb,mathtools}
\usepackage{xurl,hyperref}
\hypersetup{colorlinks=true,urlcolor=blue,linkcolor=blue}
\setcounter{secnumdepth}{0}
\setlength{\parindent}{0pt}
\setlength{\parskip}{5pt}
\setlength{\emergencystretch}{3em}
\newcommand{\sref}[2]{\hyperlink{source:#1:#2}{[S#2]}}
\newcommand{\cataloguescope}[1]{\paragraph{Source question.}#1}
\newcommand{\floor}[1]{\left\lfloor#1\right\rfloor}
\begin{document}
% UnsolvedMath ID 30000347; source catalogue code OWR-1111-001
\setcounter{section}{30000346}
\section{Three-Terminal Distance-Interdiction Complexity}\label{30000347}
{\small\sffamily UnsolvedMath ID 30000347 · Computer Science}
\subsection{Definitions and mathematical statement}

\paragraph{Shared notation.}Standard number systems are denoted by $\mathbb N,\mathbb Z,\mathbb Q,\mathbb R,\mathbb C$; index ranges, computational models and specialized notation follow the cited problem source and the definitions below.


$d(u,v)$ is shortest-path length in the original graph; after deletion it is computed in the remaining graph, with disconnected pairs assigned infinite distance. In the unweighted-cost variant, deletion cost is the number of removed edges.
\paragraph{Statement recorded in the supplied source.}
Given an undirected graph, with or without edge costs, and three vertices $r,s,t$, determine the complexity of finding a minimum-cost edge set whose deletion increases each of $d(r,s)$, $d(r,t)$, and $d(s,t)$ by at least $1$. \sref{30000347}{2}
\paragraph{Scope qualification.}
The supplied graph is undirected; the separately mentioned directed variant is not added to its target. Edge lengths and deletion costs should not be conflated.
\paragraph{Source status.}
The supplied catalogue labels this record open. Its source-status metadata was checked against the supplied snapshot on 2026-10-11. This import does not independently establish current open status or mathematical correctness. The catalogue literature assessment is dated 2026-08-21.
\subsection{Short English statement}
Classify minimum-cost edge deletion that increases all three pairwise terminal distances.
\subsection{Sources}
\begin{itemize}
\item[\hypertarget{source:30000347:1}{S1}] Ulam AI and the UnsolvedMath Contributors, UnsolvedMath: A Curated Collection of Open Mathematics Problems, record 30000347 (OWR-1111-001). Catalogue curation and metadata: CC BY 4.0. The supplied source question is retained; definitions and the English summary are editorial additions.
\url{https://huggingface.co/datasets/ulamai/UnsolvedMath}.
\item[\hypertarget{source:30000347:2}{S2}] Combinatorial Optimization (2006), Oberwolfach Reports, DOI 10.4171/owr/2005/50, pp. 2833--2932
\url{https://doi.org/10.4171/owr/2005/50}.
\textit{Underlying Oberwolfach Reports material retains its source-specific licence, including CC BY-SA 4.0 where applicable; see THIRD\_PARTY\_NOTICES.md.}
\end{itemize}
\paragraph{Preparation and review.}
Original question attributed to the cited source and UnsolvedMath contributors. Codex (GPT-6; exact serving revision not exposed) assisted with definitions, English summaries, source inspection and formatting. No new human mathematical authorship or independent proof verification is asserted. Canonical ID 30000347 and source code OWR-1111-001 checked against the complete supplied registry and existing repository statements; known flagged repetitions and obvious identical conjectures excluded. Full mathematical-equivalence/literature review has not been performed. Statement preparation and final review followed contribution prompts C and D; automated checks verify structure and identity, not open status.
\label{end:30000347}
\end{document}
